\chapter{Zakończenie}
\label{ch:zakonczenie}

\section{Wnioski z przeprowadzonych badań}
\label{sec:wnioski}

% TODO: Zsyntetyzować najważniejsze obserwacje z implementacji i eksperymentów.
% Wnioski powinny wynikać z danych przedstawionych w rozdziale 4.

\section{Realizacja celu i zakresu pracy}
\label{sec:realizacja-celu}

% TODO: Odnieść się bezpośrednio do celu i zakresu zadeklarowanych we wstępie.
% Należy porównać planowane zadania z faktycznie uzyskanymi rezultatami.

\section{Ograniczenia techniczne zaproponowanego rozwiązania}
\label{sec:ograniczenia-techniczne}

% TODO: Opisać ograniczenia: jakość przy ultra-niskiej przepływności, koszt treningu,
% zależność od danych, stabilność RVQ, opóźnienie, wymagania sprzętowe i gotowość do inferencji brzegowej.

\section{Kierunki dalszego rozwoju}
\label{sec:kierunki-rozwoju}

% TODO: Wskazać dalsze prace: lepsze kodowanie entropijne, wariant przyczynowy/streamingowy,
% optymalizację pod urządzenia brzegowe, testy subiektywne, większy zbiór danych i porównanie z kodekami referencyjnymi.
